% Generated draft: requires mathematical and source review.
\documentclass[11pt]{article}
\usepackage[margin=25mm]{geometry}
\usepackage{fontspec}
\setmainfont{Latin Modern Roman}
\usepackage{amsmath,amssymb,mathtools,unicode-math}
\setmathfont{Latin Modern Math}
\usepackage{enumitem,needspace,xurl,hyperref}
\setlength{\parindent}{0pt}
\setlength{\parskip}{5pt}
\setlength{\emergencystretch}{3em}
\newcommand{\N}{\mathbb{N}}
\newcommand{\Z}{\mathbb{Z}}
\newcommand{\Q}{\mathbb{Q}}
\newcommand{\R}{\mathbb{R}}
\newcommand{\C}{\mathbb{C}}
\newcommand{\F}{\mathbb{F}}
\newcommand{\A}{\mathbb{A}}
\newcommand{\PP}{\mathbb{P}}
\newcommand{\E}{\mathbb{E}}
\newcommand{\eps}{\varepsilon}
\newcommand{\dd}{\,\mathrm{d}}
\begin{document}
\section*{UnsolvedMath ID: 800887}
\textbf{Source problem code:} catalogue entry
\subsection{Definitions and mathematical statement}
The supplied excerpt does not explicitly specify the ambient manifold, the volume form, or the solution convention. Under the customary compact K\"ahler interpretation, let \(X\) be a compact K\"ahler manifold of complex dimension \(n\geq 1\), let \(dV\) be a smooth positive volume form, and let \(\omega\) be a smooth closed real \((1,1)\)-form. These ambient assumptions require confirmation against the full source. Assume that \(\omega\) is semipositive, meaning \(\omega_x(v,\overline v)\geq 0\) in the Hermitian interpretation, and big in the sense explicitly used in the supplied record:
\[
\int_X\omega^n>0.
\]
An \(\omega\)-plurisubharmonic function \(\varphi\) is an upper semicontinuous, locally integrable real-valued function for which \(\omega+dd^c\varphi\) is a positive current. Here \(dd^c\) is the real complex Hessian operator, with one fixed conventional normalization used throughout.

For \(p>1\), let \(f\in L^p(X,dV)\) be nonnegative and satisfy the necessary mass compatibility condition
\[
\int_X f\,dV=\int_X\omega^n.
\]
Consider a continuous \(\omega\)-plurisubharmonic solution \(\varphi\) of
\[
(\omega+dd^c\varphi)^n=f\,dV,
\]
where the Monge--Amp\`ere measure is interpreted in the Bedford--Taylor sense for locally bounded potentials. Nonnegativity, mass compatibility, and this solution convention are implicit in this reconstruction, rather than explicitly stated hypotheses of the supplied excerpt.

The question attributed to Eyssidieux--Guedj--Zeriahi is whether every such solution is H\"older continuous on all of \(X\): for a fixed smooth Riemannian distance \(d\) on \(X\), must there exist \(\alpha\in(0,1]\) and \(C<\infty\) such that
\[
|\varphi(x)-\varphi(y)|\leq C\,d(x,y)^\alpha
\qquad\text{for all }x,y\in X?
\]
No uniform exponent or quantitative dependence of \(\alpha\) and \(C\) is specified in the supplied question. Adding a constant to \(\varphi\) does not affect this property.

The supplied background attributes continuity for these \(L^p\) equations, obtained using viscosity techniques, to the reference labeled \([\mathrm{EGZ11}]\), and states uniqueness up to an additive constant. The ambient and connectedness hypotheses for that uniqueness assertion are not supplied. Its description of further regularity as largely open, and the registry's open status, are supplied source assessments rather than independently verified current literature conclusions.
\subsection{Short English statement}
The question asks whether solutions of the complex Monge--Amp\`ere equation with an \(L^p\) density, \(p>1\), are H\"older continuous everywhere when the background form is semipositive and has positive top self-intersection. The supplied background asserts continuity and uniqueness up to an additive constant. The precise ambient hypotheses and solution convention need confirmation from the full source.
\subsection{Sources}
\begin{itemize}
\item Local UnsolvedMath registry, supplied complete record: permanent ID \(800887\), exact source problem code \texttt{catalogue entry}, title \emph{Open Problems in Pluripotential Theory --- Question 18}. The supplied registry labels the problem open; this label has not been independently verified.
\item Dinew, Guedj, Zeriahi, \emph{Open Problems in Pluripotential Theory} (2015), Question 18. Supplied source URL: \url{https://arxiv.org/abs/1511.00705}. The question is explicitly attributed in the excerpt to Eyssidieux--Guedj--Zeriahi. The supplied provenance records access on 2026-10-16 and extraction by \texttt{arxiv\_latex}; these are supplied provenance claims, not a claim of access or verification for this draft.
\item Reference label \([\mathrm{EGZ11}]\) in the supplied excerpt: cited there for continuity using viscosity techniques. No full bibliographic entry or separate URL for this reference is supplied.
\end{itemize}
\end{document}
